% Part: incompleteness
% Chapter: theories-computability
% Section: interpretability

\documentclass[../../../include/open-logic-section]{subfiles}

\begin{document}

\olfileid{inc}{tcp}{itp}

\olsection{نظریه‌هایی که~$\Th{Q}$ در آن‌ها تفسیرپذیر است تصمیم‌ناپذیرند}

می‌توانیم این نتایج را باز هم تقویت کنیم. به‌طور غیرصوری، تفسیر یک
زبان~$\Lang{L_1}$ در زبان دیگری چون~$\Lang{L_2}$ مستلزم تعریف جهان،
نمادهای رابطه و نمادهای تابعِ~$\Lang{L_1}$ با !!{formula}sی در
$\Lang{L_2}$ است. گرچه در اینجا فرصت پرداختن به آن را نداریم، می‌توان
این تعریف را دقیق کرد.

\begin{thm}
  فرض کنید~$\Th{T}$ نظریه‌ای در زبانی باشد که بتوان زبان حساب را در آن
  چنان تفسیر کرد که~$\Th{T}$ با تفسیر~$\Th{Q}$ سازگار باشد. آنگاه
  $\Th{T}$ تصمیم‌ناپذیر است. اگر~$\Th{T}$ تفسیر اصول موضوعهٔ~$\Th{Q}$
  را ثابت کند، هیچ گسترش سازگاری از~$\Th{T}$ تصمیم‌پذیر نیست.
\end{thm}

اثبات تنها اصلاح کوچکی در اثبات قضیهٔ پیشین است؛ از یک مثال نقض می‌توان
برای به‌دست‌آوردن جداسازی~$\Th{Q}$ و~$\Th{\bar Q}$ استفاده کرد.
می‌توان~$\Th{ZFC}$، یعنی نظریهٔ مجموعه‌های زرملو--فرانکل همراه با اصل
انتخاب، را بنیادی اصل‌موضوعی دانست که برای بخش بزرگی از ریاضیات معمول
توان کافی دارد. در~$\Th{ZFC}$ می‌توان اعداد طبیعی را تعریف کرد و در
چارچوب این تفسیر، اصول موضوعهٔ~$\Th{Q}$ صادق‌اند. بنابراین داریم:

\begin{cor}
هیچ گسترش سازگار و تصمیم‌پذیری از~$\Th{ZFC}$ وجود ندارد.
\end{cor}

\begin{cor}
هیچ گسترش کامل و سازگاری از~$\Th{ZFC}$ وجود ندارد که به‌طور
محاسبه‌پذیر !!{axiomatizable} باشد.
\end{cor}

زبان~$\Th{ZFC}$ تنها یک رابطهٔ دودویی، یعنی~$\in$، دارد. (در واقع حتی
به همانی هم نیاز ندارید.) بنابراین داریم:

\begin{cor}
منطق مرتبه‌اول برای هر زبانی که نماد رابطه‌ای دودویی داشته باشد
تصمیم‌ناپذیر است.
\end{cor}

این نتیجه به هر زبانی با دو نماد تابع یک‌موضعی گسترش می‌یابد، زیرا
می‌توان از آن‌ها برای شبیه‌سازی نماد رابطه‌ای دودویی استفاده کرد. نتایج
یادشده دقیق‌اند: معلوم می‌شود منطق مرتبه‌اول برای زبانی که فقط نمادهای
رابطهٔ \emph{یک‌موضعی} و حداکثر یک نماد تابع \emph{یک‌موضعی} دارد
تصمیم‌پذیر است.

نکتهٔ پایانی: می‌دانیم مجموعهٔ جمله‌های زبان~$\Obj 0$، $'$، $+$،
$\times$، $<$ که در مدل استاندارد صادق‌اند تصمیم‌ناپذیر است. در واقع
می‌توان~$<$ را با نمادهای دیگر تعریف کرد، و سپس~$+$ را با~$\times$
و~$'$ تعریف کرد. پس مجموعهٔ جمله‌های صادق در زبان~$\Obj 0$، $'$،
$\times$ تصمیم‌ناپذیر است. از سوی دیگر، پرسبرگر نشان داده است مجموعهٔ
جمله‌های زبان~$\Obj 0$، $'$، $+$ که در مدل استاندارد صادق‌اند تصمیم‌پذیر
است. بااین‌حال، این روش از نظر محاسباتی شدنی نیست.

\end{document}
